\documentclass[aspectratio=169,10pt]{beamer}
\usetheme[numbering=fraction,progressbar=frametitle]{metropolis}
\usepackage{fontspec}
\setsansfont{Latin Modern Sans}
\usepackage{amsmath,amssymb,mathtools,graphicx,booktabs}
\usepackage{polyglossia}
\setdefaultlanguage{french}
\setbeamertemplate{navigation symbols}{}
\setbeamerfont{frametitle}{size=\large}
\setbeamerfont{normal text}{size=\normalsize}
\setlength{\parskip}{.45em}
\hypersetup{pdfauthor={},pdfsubject={Lecture slides},pdfkeywords={},pdftitle={Calcul dans \ensuremath{\mathbb{R}}}}
\providecommand{\N}{\mathbb{N}}
\providecommand{\Z}{\mathbb{Z}}
\providecommand{\D}{\mathbb{D}}
\providecommand{\Q}{\mathbb{Q}}
\providecommand{\R}{\mathbb{R}}
\providecommand{\C}{\mathbb{C}}
\providecommand{\vect}[1]{\overrightarrow{#1}}
\providecommand{\vi}{\vec{\imath}}
\providecommand{\vj}{\vec{\jmath}}
\providecommand{\vk}{\vec{k}}
\providecommand{\norm}[1]{\left\lVert #1\right\rVert}
\providecommand{\vc}[2]{\begin{pmatrix}#1\\#2\end{pmatrix}}
\providecommand{\deter}[2]{\det\left(#1,#2\right)}
\providecommand{\abs}[1]{\left\lvert #1\right\rvert}
\providecommand{\dd}{\,\mathrm{d}}
\providecommand{\e}{\mathrm{e}}
\providecommand{\Card}{\operatorname{Card}}
\providecommand{\E}{\operatorname{E}}
\providecommand{\V}{\operatorname{V}}
\providecommand{\degre}{^{\circ}}
\providecommand{\un}[1]{\,\mathrm{#1}}
\providecommand{\sci}[2]{#1\times10^{#2}}
\providecommand{\pgcd}{\operatorname{PGCD}}
\providecommand{\ppcm}{\operatorname{PPCM}}
\title{Calcul dans \ensuremath{\mathbb{R}}}
\subtitle{Cours et exemples}
\date{}
\author{}
\begin{document}
\begin{frame}[plain]
\vfill
{\usebeamerfont{title}\inserttitle\par}
\medskip
{\usebeamerfont{subtitle}\insertsubtitle\par}
\medskip
{\usebeamercolor[fg]{progress bar}\rule{\textwidth}{.5pt}}
\vfill
\end{frame}
\begin{frame}{}
Distribuer des cartes portant des nombres de formes différentes. Construire un
diagramme d'inclusions pour \(\mathbb{N}\), \(\mathbb{Z}\), \(\mathbb{D}\), \(\mathbb{Q}\) et \(\mathbb{R}\), puis placer chaque
carte dans le plus petit ensemble possible. Discuter les cartes \(0{,}333\ldots\),
\(\sqrt{9}\) et \(\sqrt2\).
\end{frame}
\begin{frame}{Objectifs}
\begin{itemize}

  \item Maîtriser la hiérarchie \(\mathbb{N} \subset \mathbb{Z} \subset \mathbb{Q} \subset \mathbb{R}\).
  
\item Calculer avec les fractions, puissances et racines carrées.
  
\item Utiliser la valeur absolue et résoudre des équations/inéquations.
  
\item Utiliser les notations d'intervalles et d'encadrements.

\end{itemize}
\end{frame}
\begin{frame}{Les ensembles de nombres}
\[\mathbb{N} \subset \mathbb{Z} \subset \mathbb{Q} \subset \mathbb{R}\]

\begin{itemize}

  \item \(\mathbb{N} = \{0, 1, 2, 3, \ldots\}\) (entiers naturels)
  
\item \(\mathbb{Z} = \{\ldots, -2, -1, 0, 1, 2, \ldots\}\) (entiers relatifs)
  
\item \(\mathbb{Q} = \bigl\{\dfrac{p}{q} \mid p \in \mathbb{Z},\ q \in \mathbb{N}^*\bigr\}\) (rationnels)
  
\item \(\mathbb{R}\) : tous les réels, incluant les irrationnels (\(\sqrt{2}\), \(\pi\), \(e\), \ldots{})

\end{itemize}
\end{frame}
\begin{frame}{Les ensembles de nombres}
Pour diviser \(157\) par \(12\), on cherche le plus grand multiple de \(12\) qui ne
dépasse pas \(157\). La disposition ci-dessous est celle attendue sur une copie :

On lit \(157=12\times13+1\) : le quotient est \(13\) et le reste est \(1\).
Toujours vérifier que \(0\leq r<12\).
\end{frame}
\begin{frame}{Puissances et racines}
Pour \(a, b \in \mathbb{R}^*\) et \(m, n \in \mathbb{Z}\) :
\[a^m \times a^n = a^{m+n},\quad \dfrac{a^m}{a^n} = a^{m-n},\quad
  (a^m)^n = a^{mn},\quad (ab)^n = a^n b^n.\]
\end{frame}
\begin{frame}{Puissances et racines}
Pour \(a \geq 0\), \(b > 0\) :
\[\sqrt{ab} = \sqrt{a}\,\sqrt{b},\qquad
  \sqrt{\dfrac{a}{b}} = \dfrac{\sqrt{a}}{\sqrt{b}},\qquad
  \bigl(\sqrt{a}\bigr)^2 = a.\]
\end{frame}
\begin{frame}{Valeur absolue}
Pour tout réel \(x\) :
\[\left|x\right| = \begin{cases} x & \text{si } x \geq 0, \\ -x & \text{si } x < 0. \end{cases}\]
Géométriquement, \(\left|x\right|\) est la \emph{distance} de \(x\) à l'origine.
\(\left|x - a\right|\) est la distance de \(x\) au point \(a\).
\end{frame}
\begin{frame}{Valeur absolue}
Pour tous réels \(a\) et \(b\) : \(\left|a + b\right| \leq \left|a\right| + \left|b\right|\).
\end{frame}
\begin{frame}{Valeur absolue}
On travaille sur le domaine de \(f\) et on suppose \(k\geq0\).
Si \(k<0\), les conditions \(|f(x)|\leq k\) et \(|f(x)|=k\) sont impossibles,
et \(|f(x)|\geq k\) est vraie sur tout le domaine.

\begin{enumerate}
\setcounter{enumi}{0}
  \item \(\left|f(x)\right| \leq k \Leftrightarrow -k \leq f(x) \leq k\).
  
\item \(\left|f(x)\right| \geq k \Leftrightarrow f(x) \leq -k \text{ ou } f(x) \geq k\).
  
\item \(\left|f(x)\right| = k \Leftrightarrow f(x) = k \text{ ou } f(x) = -k\).

\end{enumerate}
\end{frame}
\begin{frame}{Intervalles et encadrements}
Un \textbf{intervalle} est un sous-ensemble de \(\mathbb{R}\) formé de tous les réels
compris entre deux bornes. On distingue :
\end{frame}
\begin{frame}{Intervalles et encadrements (2)}
\begin{center}\begin{tabular}{lll}
Notation & Définition & Représentation \\ 
\([a,b]\) & \(a \leq x \leq b\) &  \\ 
\(]a,b[\) & \(a < x < b\) &  \\ 
\([a,b[\) & \(a \leq x < b\) &  \\ 
\(]a,b]\) & \(a < x \leq b\) &  \\ 
\([a,+\infty[\) & \(x \geq a\) &  \\ 
\(]-\infty,b]\) & \(x \leq b\) &  \\ 
\(\mathbb{R} = ]-\infty,+\infty[\) & tous les réels &  \\ 
\end{tabular}\end{center}
\end{frame}
\begin{frame}{Intervalles et encadrements}
\(+\infty\) et \(-\infty\) ne sont \textbf{pas} des réels : on écrit toujours
\(]a,+\infty[\) et \(]-\infty,b[\) (crochets ouverts).
\end{frame}
\begin{frame}{Intervalles et encadrements}
\begin{itemize}

  \item \([1,4] \cap [3,7] = [3,4]\) (valeurs \emph{communes} aux deux intervalles)
  
\item \([1,4] \cup [3,7] = [1,7]\) (valeurs dans \emph{au moins un} des deux)
  
\item \([1,3] \cap [5,7] = \emptyset\) (intervalles disjoints)

\end{itemize}
\end{frame}
\begin{frame}{Intervalles et encadrements}
Pour encadrer \(\sqrt{n}\) entre deux entiers consécutifs :

\begin{enumerate}
\setcounter{enumi}{0}
  \item Trouver \(k\) tel que \(k^2 \leq n < (k+1)^2\).
  
\item Conclure : \(k \leq \sqrt{n} < k+1\).

\end{enumerate}

\emph{Exemple :} \(4^2=16 < 17 < 25=5^2\), donc \(4 < \sqrt{17} < 5\).
\end{frame}
\begin{frame}{Fractions et calcul algébrique}
Pour \(a,c\in\mathbb{R}\) et \(b,d\in\mathbb{R}^*\) (avec aussi \(c\ne0\) pour le quotient) :
\[\dfrac{a}{b} + \dfrac{c}{d} = \dfrac{ad+bc}{bd},\qquad
  \dfrac{a}{b} \times \dfrac{c}{d} = \dfrac{ac}{bd},\qquad
  \dfrac{a/b}{c/d} = \dfrac{a}{b} \times \dfrac{d}{c} = \dfrac{ad}{bc}.\]
\end{frame}
\begin{frame}{Fractions et calcul algébrique}
Pour éliminer une racine carrée au dénominateur, multiplier numérateur
et dénominateur par l'\emph{expression conjuguée} non nulle.
Pour \(a,b\geq0\) et \(a\ne b\) :
\[\dfrac{1}{\sqrt{a}+\sqrt{b}} \times \dfrac{\sqrt{a}-\sqrt{b}}{\sqrt{a}-\sqrt{b}}
  = \dfrac{\sqrt{a}-\sqrt{b}}{a-b}.\]
Si \(a=b>0\), on utilise directement \(1/(2\sqrt a)=\sqrt a/(2a)\).
Si \(a=b=0\), le quotient initial n'est pas défini.
\end{frame}
\begin{frame}{Fractions et calcul algébrique}
\begin{itemize}

  \item \(\dfrac{1}{\sqrt{3}} = \dfrac{\sqrt{3}}{3}\)  
        (multiplier par \(\sqrt{3}/\sqrt{3}\))
  
\item \(\dfrac{2}{\sqrt{5}-1} = \dfrac{2(\sqrt{5}+1)}{(\sqrt{5})^2-1^2}
        = \dfrac{2(\sqrt{5}+1)}{4} = \dfrac{\sqrt{5}+1}{2}\)

\end{itemize}
\end{frame}
\end{document}
